\section{Limitations}
\label{sec:platform-limitations}

The platform is a first version, and its boundaries should be stated.
\begin{itemize}
  \item \emph{Comparable shape, not comparable cost.} The trajectory contract
    aligns methods on iteration count. One model-free step, one classical
    iteration and one unrolled layer do very different amounts of work, and
    the harness does not measure FLOPs or wall-clock per step. Any
    matched-budget statement is iteration-matched, which favours the
    model-free optimizers.
  \item \emph{Coverage.} Only LISTA and ALISTA of the nine unrolled variants
    in the model-based library are wired into the harness, and only they
    consume the shared $x_0$. Classical-optimizer support depends on the host
    library: the L2O-Scale family exposes SGD, Adam and Adagrad, the others
    only Adam. The L2O-Scale family trains and evaluates on different problem
    sets by design, which the adapter maps separately.
  \item \emph{Unstandardised per-instance breakdowns.} Beyond the four
    trajectory summaries, the optional \texttt{per\_seed} field of an
    evaluation result has an adapter-specific shape (per seed, per test file,
    or otherwise), so generic downstream code cannot interpret it without
    knowing the adapter.
  \item \emph{Suites are not problem-checked.} Nothing prevents a suite from
    ranking methods evaluated on different problems in one leaderboard; suites
    must be kept problem-scoped by the user.
  \item \emph{Evaluation determinism} is exact for the file-based LASSO
    problems but not for the neural-network optimizees of the L2O-Scale
    family (Section~\ref{sec:platform-repro}).
  \item \emph{Scale.} All results were produced on one 8\,GB GPU inside one
    container; the platform has no notion of distributing a sweep across
    machines, and sweeps are hours to days long.
\end{itemize}
